\section{Policy evaluation software must preserve the decision and its uncertainty}
\label{exec23:scope-rl-scope}
A servicing team may use an offline policy evaluation library to compare a
sequence of reminders, forbearance offers and escalations. The useful output is
the value of a specified feasible policy under an information clock, together
with defensible uncertainty. A function returning a ranking is only one part of
that task. SCOPE-RL provides a particularly useful proposal comparator because
its selected discrete estimators, interval helpers and policy-selection paths
make several boundaries explicit. This review concerns version 0.2.1 at the
immutable 1 December 2023 revision returned by the cutoff query. Eleven
authenticated source and declaration files support the following observations;
the continuous-action, distributional and density-learning families have not
received equivalent inspection. No upstream program was executed.
\citep{exec23_scope_rl_sources}

\subsection{An episode, a propensity and a policy are substantive inputs}
\label{exec23:scope-rl-target}
Let a customer episode have $T$ decisions. At decision $t$, $H_t$ is all
information available before the action, $A_t$ the action, and $R_t$ the
subsequent monetary reward. Write $b_t(a\mid H_t)$ for the actual logging
probability and $\pi_t(a\mid H_t)$ for the policy being evaluated. For a declared
initial population and discount $0<\gamma\leq1$, the proposed target is
\[
 J(\pi)=\mathbb E_\pi\!\left[\sum_{t=0}^{T-1}\gamma^tR_t\right],
 \qquad
 W_{it}=\prod_{s=0}^{t}
 \frac{\pi_s(A_{is}\mid H_{is})}{b_s(A_{is}\mid H_{is})}.
\]
The elementary change-of-measure argument needs the logging probabilities to
describe the actual action mechanism, support for every target action-history
pair, and the same conditional transition/reward law under the intervention.
Multiplying the logging trajectory density by $W_{i,T-1}$ then replaces each
action factor by its target counterpart. For a reward observed by time $t$,
future action factors integrate to one, which explains the per-decision form.
An omitted private distress signal or a post-action feature can invalidate this
argument even if every probability row sums to one. The independent-episode
sampling condition used for the simple uncertainty examples below is additional
to this identification argument.

The inspected implementation reshapes flattened records into equal-length
episodes, cumulatively multiplies action probabilities, and either repeats the
final weight or retains the prefix weights. Its two scores therefore correspond
to
\[
 Z_i^{\rm TIS}=W_{i,T-1}\sum_t\gamma^tR_{it},
 \qquad Z_i^{\rm PDIS}=\sum_t\gamma^tW_{it}R_{it}.
\]
The interface checks dimensions, row counts, divisibility by the horizon and
target probability normalization. Such checks cannot establish truthful
logging, temporal ordering or common episode membership. The probability
validator is called with a zero lower limit; the discount expression divides a
cumulative product by $\gamma$ while its scalar check also uses a zero lower
limit. A proposed adapter must enforce finite, strictly positive observed-action
denominators and an explicitly supported discount domain before these calls.
This is a static source finding, not an observed runtime failure.
\citep{exec23_scope_rl_sources}

An exact two-period example clarifies why two estimators can disagree without
targeting different policies. Let both logging actions be independent fair
coins and both target actions have probability $3/4$ of action one. Set
$R_0=2A_0$, $R_1=4+4A_1$ and $\gamma=1/2$. Direct enumeration of all four
trajectories gives $J(\pi)=5$. Both importance scores average to five under the
logging law. On the particular two-episode sample $(1,0),(0,1)$, however, TIS is
$3$ and PDIS is $15/4$. Neither sample result proves bias or superior precision.
The project's rational-arithmetic reference enumerates the complete finite law
and distinguishes this population calculation from the illustrative sample.
\citep{exec23_scope_rl_reference}

The inspected doubly robust score adds a fitted state-value term and subtracts
the fitted action value, with current and previous prefix weights:
\[
 Z_i^{\rm DR}=\sum_t\gamma^t
 \{W_{it}(R_{it}-\widehat Q_t(H_{it},A_{it}))
       +W_{i,t-1}\widehat V_t(H_{it})\},\quad W_{i,-1}=1.
\]
Its source calculation is useful for conversion tests, but the label does not
prove robustness to an arbitrary hidden logging signal, an estimated policy,
or dependent customers. The input builder defaults to one fitted-Q fold; its
multi-fold branch excludes contiguous groups of whole trajectories. It caches
fitted models by policy name. Therefore episode grouping, calendar dependence,
policy identity and changes to training settings must be bound explicitly.
Repeated calls with a reused name are not evidence of a fresh fit. The d3rlpy
delegate and a full frozen environment remain native qualification obligations.
\citep{exec23_scope_rl_sources}

\subsection{Terminal rewards and uncertainty cannot be inferred from names}
\label{exec23:scope-rl-uncertainty}
The high-level evaluator defaults to disabling rewards after a done flag. The
actual mask is the cumulative product of $1-d_t$, including the current row.
If a dataset defines $d_t$ as indicating that the transition just produced the
terminal reward, this removes that terminal reward as well. A different
pre-transition convention could justify the mask; the adapter must state and
test the convention. With undiscounted rewards $(2,8)$ and flags $(0,1)$, the
literal mask leaves value $2$, whereas retaining the terminal receipt gives
$10$. At a stipulated servicing cost of six dollars, the margins are $-4$ and
$4$. These are authored examples, not measured bank returns or observed
package outputs. The dataset documentation separately describes episode end
and maximum-step termination, while the fitted-Q dataset call passes the done
array as its terminal argument. A conversion must resolve both meanings.
\citep{exec23_scope_rl_sources,exec23_scope_rl_reference}

The basic estimator's interval method first calculates one score per trajectory
and delegates the resulting vector to a common interval helper. The bootstrap
helper samples those fixed scores with replacement, takes percentile endpoints,
and reports the mean of bootstrap means. Its signature defaults to 100
replications although its parameter prose says 10,000, and it requires an
explicit random seed. Neither this helper nor the selected call refits the
value model, estimates new logging propensities or repeats policy learning in
each bootstrap sample. Whether conditioning on those fits is scientifically
appropriate needs a separate argument; it is not supplied by the word
bootstrap. \citep{exec23_scope_rl_sources}

This distinction already matters without a learned nuisance. For a
self-normalized weighted mean, let two episode weights be $(1,3)$ and returns
$(0,4)$. Ignoring the implementation's tiny denominator stabilizer solely for
this exact algebraic illustration, the estimate is $3$ and the precomputed
normalized scores are $(0,6)$. The four equally likely size-two bootstrap index
pairs produce means $(0,3,3,6)$ if these scores are fixed, but ratios
$(0,3,3,4)$ if the denominator is recalculated. Their bootstrap variances are
$9/2$ and $9/4$. This is not a coverage assessment of either procedure; it shows
that they are different resampling targets. The selected self-normalized class
inherits the score-based interval path and uses a $10^{-10}$ denominator
stabilizer. A native test must retain that stabilizer, distinguish ratio
recalculation from score resampling, and state which procedure is justified.
\citep{exec23_scope_rl_sources,exec23_scope_rl_reference}

Three analytic interval helpers need particularly careful admission. The
Hoeffding-named helper computes the radius
\[
 \max_i z_i\sqrt{\log(2/\alpha)n/2},
\]
whereas its adjacent displayed formula places $n$ in the denominator. Parsing
the expression as inert syntax confirms multiplication by $n$; at $n=20$ the
source radius is twenty times the expression with the same logarithm and $n$
in the denominator. More fundamentally, an observed sample maximum is not a
known population range. For independent Bernoulli rewards with mean $1/100$,
all twenty observations are zero with probability $(99/100)^{20}>0.8$. The
source Hoeffding and Bernstein expressions then both give zero radius despite
the positive mean. A constant negative sample also makes the first radius
negative and reverses its endpoints. No concentration guarantee follows from
these helper names for unrestricted banking rewards.
\citep{exec23_scope_rl_sources,exec23_scope_rl_reference}

The Bernstein-named helper also differs from its printed variance expression:
the executable expression divides the variance contribution by $n-1$, while
the prose display includes $n(n-1)$, and a supplied sample maximum again replaces
a declared range. These observations warrant explicit comparison of definitions
and sample-variance conventions, not a silent formula patch. The Student helper
uses the $1-\alpha$ quantile and returns symmetric endpoints labelled
$100(1-\alpha)$ percent. Under the exact nondegenerate normal-sample model in
which the usual Student pivot has a symmetric continuous law, those endpoints
have central probability $1-2\alpha$: $90$ percent at $\alpha=0.05$. This
algebra does not establish a distributional approximation for importance
scores. Positive sample sizes, finite values, signed rewards, degenerate
samples, alpha endpoints and valid simultaneous coverage must all be explicit
adapter contracts. \citep{exec23_scope_rl_sources,exec23_scope_rl_reference}

\subsection{A selection benchmark is not an offline safety certificate}
\label{exec23:scope-rl-decision}
The inspected lower-bound selector obtains on-policy selection results before
ranking candidate interval endpoints, even when the requested output omits
evaluation metrics. The called routine rejects a missing on-policy return
array. This is a concrete benchmark-data requirement: a logged-only banking
request generally cannot furnish those counterfactual returns. Missing truth
must lead to a refusal or a separately designed selection path, never fabricated
on-policy values. The selected call forwards the same alpha to each candidate
and estimator; no multiplicity adjustment appears in this path. Even twenty
independent tests with individual false-admission probability $0.05$ have
probability $1-0.95^{20}$, approximately $0.642$, of at least one false admission.
That hypothetical example is not a measured SCOPE-RL error rate. A candidate
policy learned or selected on the evaluation histories requires its own
selection-aware design. \citep{exec23_scope_rl_sources,exec23_scope_rl_reference}

The ground-truth lower-quartile output also needs a name/value fixture. The
source sorts the policy-value array before deriving a lower-quartile name
ordering, then sorts lower-quartile values separately. In input order A, B,
give A returns $(-10,-10,40,40)$ and B returns $(4,4,4,4)$. Their means are
$15,4$, but lower quartiles are $-10,4$. The selected expressions pair the
unchanged name order A, B with descending values $4,-10$, rather than the
correct lower-tail name order B, A. The static witness is sufficient to require
an identity-preserving test; native behavior in a frozen dependency environment
has not been measured. A printed safety or regret column must also preserve the
source's explicit true-safe and estimated-safe denominators instead of relying
on generic type-I/type-II labels. \citep{exec23_scope_rl_sources,exec23_scope_rl_reference}

The proposal disposition is therefore \emph{reference}, with no approved native
adapter. SCOPE-RL's explicit trajectory scores, fitted-value preparation and
evaluation organization inform a narrow design. Whole-package adoption would
admit much more functionality than the inspected contract. A Python worker is
a candidate only after complete artifact and dependency custody, terminal/time
fixtures, interval repair review and fresh native conformance. Porting the
package to R, Java or C++ currently has no demonstrated advantage sufficient to
justify reproducing its preprocessing, model delegates and maintenance burden.
Independent exact references are valuable test oracles; they do not replace the
package or establish empirical policy performance.

The source declares Apache 2.0 and includes the license text, but this does not
clear transitive dependencies, datasets or deployment terms. The requirements
span NumPy, SciPy, pandas, scikit-learn, plotting libraries, torch, d3rlpy,
Gym/Gymnasium and Hydra; most are lower bounds, not a reproducible environment
lock. Named authors and the dated revision establish identity, not a current
support commitment. The maintenance owner must resolve actual dependency
versions and rights; the scientific owner must define target, logging truth,
support and sampling law; the operator must measure resource and recovery cost.
\citep{exec23_scope_rl_sources}

A proposed library/CLI/MCP request should carry immutable episode IDs, decision
times, the information set, reward/termination conventions, action vocabulary,
logging and target probabilities, policy/training identities, horizon, discount,
uncertainty method and random seed. Its response should preserve the target,
score definition, source revision, refusal reason, finite-weight diagnostics and
whether uncertainty is conditional on fitted quantities. An agent must not
convert a schema-valid answer into an identified effect or a safety approval.
Public qualification can begin with finite known-policy trajectory laws and
adversarial permutations, missing propensities, terminal rewards, signed returns,
shared calendar shocks and policy-selection leakage. Fresh native tests and
measured costs follow only in a reviewed environment; bank truth and adoption
remain later independent gates.
